% !TeX root = bash.tex
\section{V. Bash Scripting}
\begin{frame}[fragile]
\frametitle{Bash is a programming language}

Besides executing commands, bash is also a programming language. We can write
and execute bash scripts.
\newline \newline
Let's say you want to print many files at once.
\newline \newline
Try this simple example:
\begin{lstlisting}[language=bash]
$ for i in {01..05}; do \
    cat um-logo-$i.txt; \
  done
\end{lstlisting}
\end{frame}

\begin{frame}[fragile]
\frametitle{Special characters}
Some characters in bash has non-literal meaning. Here's a list of them:
\cverb+ $'"\#=[]!><|;{}()*?~\&+. Use backslash \cverb|\| to escape special
characters.
\newline \newline
Whitespace separates words. To prevent word separation, use single or double
quotes:
\begin{lstlisting}[language=bash]
$ mkdir milias lab reports    # Creates three folders
$ mkdir 'milias lab reports'  # Creates one folder
$ mkdir "milias lab reports"  # Also works
\end{lstlisting}
Everything inside single quotes has literal meaning. You can't use single
qoutes inside single quotes; use double quotes in such case. More about double
qoutes will be covered later.
\end{frame}

\begin{frame}[fragile]
\frametitle{Variables}
\textbf{Variables} in bash are generally strings.
A variable is defined this way:
\begin{lstlisting}[language=bash]
$ i=0
$ s=s
\end{lstlisting}
Surprisingly these do \textbf{not} work:
\begin{lstlisting}[language=bash]
$ i = 0  # WRONG
$ s = s  # WRONG
\end{lstlisting}
They can be accessed with a dollar sign:
\begin{lstlisting}[language=bash]
$ echo $i  # 0 P.S. This is a string
$ echo $s  # s
\end{lstlisting}
\end{frame}

\begin{frame}[fragile]
\frametitle{Variable expansion}
The use of the dollar sign in the previous example is called
\textbf{variable expansion}. Variable expansion works inside double quotes, but
not single quotes. You may enclose variable names in braces when
necessary:
\begin{lstlisting}[language=bash]
$ fruit=apple
$ echo "The plural of $fruit is $fruits."   # Don't do this
$ echo "The plural of $fruit is ${fruit}s." # Do this
\end{lstlisting}
\end{frame}

\begin{frame}[fragile]
\frametitle{Environment variables}
An \textbf{environment variable} can be set and accessed like this:
\begin{lstlisting}[language=bash]
$ export VAR=VALUE  # set
$ echo $VAR         # access
\end{lstlisting}
\cverb|export| exposes the variable to child processes.
\end{frame}

\begin{frame}[fragile]
\frametitle{Environment variables}
Suppose you're downloading something from a completely legal website,
and you want the traffic to go through your completely legal local proxy
for completely legal reasons.
\begin{lstlisting}[language=bash]
# dont actually execute this unless you have a proxy listening
# on port 8080
# set environment variable for local proxy
$ export HTTPS_PROXY=http://localhost:8080/
# download the thing
$ curl -O https://legal.website/legal-thing
\end{lstlisting}
Here \cverb|export| is necessary because we need to expose the variable
\cverb|HTTPS_PROXY| to \cverb|curl|, a child process.
\end{frame}

\begin{frame}[fragile]
\frametitle{Shell scripts}
Open your favorite text editor and edit \tt{05-scripting/um.sh}:
\begin{lstlisting}[language=bash]
# 05-scripting/um.sh
for i in {01..05}; do
    cat um-logo-$i.txt
done
\end{lstlisting}
Save file, then come back to bash, \cverb|cd| into \tt{05-scripting} and run:
\begin{lstlisting}[language=bash]
$ ./um.sh
\end{lstlisting}
\end{frame}

\begin{frame}[fragile]
\frametitle{Viewing Permissions: \texttt{ls -l}}
The command didn't work because we didn't ask for permissions.
Before changing permissions, we need to know how to read them. 
We use the \texttt{ls -l} (long format) command.

\begin{lstlisting}[language=bash]
$ ls -l um.sh
-rwxr-xr-- 1 ubuntu devgroup 1048 Mar 21 10:00 um.sh
\end{lstlisting}

\begin{block}{Anatomy of the output}
The most important part is the very first column: \textbf{\texttt{-rwxr-xr--}}
\begin{itemize}
    \item It consists of exactly \textbf{10 characters}.
    \item It tells us the file type and who can do what.
    \item The rest of the line shows: links, \textbf{owner} (ubuntu), \textbf{group} (devgroup), size, date, and filename.
\end{itemize}
\end{block}
\end{frame}

\begin{frame}[fragile]
\frametitle{Decoding the 10 Characters}
Let's break down \textbf{\texttt{-rwxr-xr--}}:

\begin{itemize}
    \item \textbf{Char 1 (Type):} \texttt{-} means regular file. 
    \item \textbf{Chars 2-4 (User/Owner):} \texttt{rwx} \\
          The owner can \textbf{R}ead, \textbf{W}rite, and e\textbf{X}ecute.
    \item \textbf{Chars 5-7 (Group):} \texttt{r-x} \\
          The group can \textbf{R}ead and e\textbf{X}ecute, but \textbf{NOT} write (\texttt{-}).
    \item \textbf{Chars 8-10 (Others):} \texttt{r--} \\
          Everyone else can only \textbf{R}ead.
\end{itemize}

\begin{table}
    \centering
    \begin{tabular}{c c c c}
        \textbf{Type} & \textbf{User} & \textbf{Group} & \textbf{Others} \\
        \texttt{-} & \texttt{r w x} & \texttt{r - x} & \texttt{r - -}
    \end{tabular}
\end{table}
\end{frame}

\begin{frame}[fragile]
\frametitle{From Letters to Numbers (\texttt{chmod})}
Now we can connect the visual output of \texttt{ls -l} to the numeric mode used in \texttt{chmod}.
\vspace{0.5em}

Remember the values: \textbf{r=4, w=2, x=1, -=0}

\begin{lstlisting}[language=bash]
# Visual permission:  -  r w x   r - x   r - x
# Math equivalent:       4+2+1   4+0+1   4+0+1
# Final Number:            7       5       5
\end{lstlisting}

So, if you see \texttt{-rwxr-xr-x}, it means the file currently has \textbf{755} permissions.
\vspace{0.5em}

\begin{block}{Quick Practice}
What is the numeric mode for \texttt{-rw-r--r--}?
\pause % This adds an animation pause in Beamer!
\newline
\textbf{Answer:} 644 (4+2+0, 4+0+0, 4+0+0). Very common for text files!
\end{block}
\end{frame}

\begin{frame}[fragile]
\frametitle{File Permissions \& \texttt{chmod}}
In Linux, every file has permissions for 3 roles: \textbf{User} (owner), \textbf{Group}, and \textbf{Others}.
\begin{itemize}
    \item \textbf{r} (read) = 4
    \item \textbf{w} (write) = 2
    \item \textbf{x} (execute) = 1
\end{itemize}
We use \cverb|chmod| to change these permissions:
\begin{lstlisting}[language=bash]
# Numeric mode (Add up the numbers)
# 7=4+2+1 (rwx), 5=4+1 (r-x)
$ chmod 755 um.sh  # Owner rwx, Group r-x, Others r-x

# Symbolic mode
$ chmod +x um.sh   # Add execute permission for all
$ chmod u+w um.sh  # Add write permission for owner (user)
\end{lstlisting}
\end{frame}

\begin{frame}[fragile]
\frametitle{File Ownership: \texttt{chown}}
Sometimes you need to change who owns a file or directory. We use \cverb|chown| (change owner).
\newline \newline
This usually requires \cverb|sudo| (superuser privileges).
\begin{lstlisting}[language=bash]
# Syntax: chown [user]:[group] file

# Change owner and group to root
$ sudo chown root:root config.txt

# Recursively (-R) change owner of a whole directory
$ sudo chown -R myuser:mygroup /var/www/
\end{lstlisting}
\end{frame}

\begin{frame}[fragile]
\frametitle{Execute vs. Source a Script}
There are two distinct ways to run a shell script:
\begin{block}{1. Execute (\ctexttt{./script.sh})}
Runs in a \textbf{new subshell} (child process).
Variables and directory changes (\cverb|cd|) inside the script will \textbf{NOT} affect your current terminal. Requires \cverb|+x| permission.
\end{block}

\begin{block}{2. Source (\ctexttt{source script.sh} or \ctexttt{. script.sh})}
Runs in the \textbf{current shell}.
Variables, \cverb|export|, and directory changes \textbf{WILL} remain in your current terminal. Does \textbf{not} require \cverb|+x| permission (only read permission).
\end{block}

\begin{lstlisting}[language=bash]
# Try writing a script with "cd ../" and test both:
$ chmod +x test_run.sh
$ ./test_run.sh      # You stay in the same directory
$ source test_run.sh # You are now in /turball
\end{lstlisting}
\end{frame}

\begin{frame}[fragile]
\frametitle{Shebang for shell scripts}
Add the following line to the top of \texttt{05-scripting/um.sh}:
\begin{lstlisting}[language=bash]
#!/bin/bash
\end{lstlisting}
Then execute the following command:
\begin{lstlisting}[language=bash]
$ chmod +x um.sh
\end{lstlisting}
Now we can execute the script with the following command:
\begin{lstlisting}[language=bash]
$ ./um.sh
\end{lstlisting}
\begin{block}{Explanation}
The additional line is called shebang, which allows us to conveninently execute
shell scripts.\footnote{More information available in
\href{https://en.wikipedia.org/wiki/Shebang_(Unix)}{this Wikipedia article}}
The \cverb|chmod| command gives necessary execute permission to
our script.
\end{block}
\end{frame}


\begin{frame}[fragile]
\frametitle{Exit code}
Every command results in an \textbf{exit code} whenever it terminates.
By convention, an exit code of 0 implies success, and everything else means
something went wrong (consult respective man pages).
\newline \newline
For this very reason, in bash, \textbf{0 is boolean true} and
\textbf{everything else is false}.
\newline \newline
The sign \cverb|!| is used to negate an exit code: \cverb|! EXPR| is true only
if \cverb|EXPR| is false. The space between \cverb|!| and \cverb|EXPR| is
required.

\begin{block}{Note}
    When you \ctexttt{return 0;} at the end of \ctexttt{int main()}, you are
    emitting an exit status of zero.
\end{block}
\end{frame}

\begin{frame}[fragile]
\frametitle{Control operators}

The operators \cverb|&&| and \cverb!||! are called \textbf{control operators}.
They are used between two commands, and the first command determines
whether the second one is executed.
\begin{lstlisting}[language=bash]
$ mkdir dir && cd dir
\end{lstlisting}
Here \cverb|cd| is executed only if \cverb|mkdir| is successful. \cverb!||!
does the exact opposite: the second command is executed only if the first one
fails.
\end{frame}

\begin{frame}[fragile]
\frametitle{\textbf{if} statements}
Usually, we use an \textbf{if} statement to:
\begin{itemize}
    \item Run a command and see if it succeeds
    \item Test the value of a variable
    \item Check if a file exists
\end{itemize}
It looks like this:
\begin{lstlisting}[language=bash]
if COMMAND; then
    BODY
elif COMMAND; then
    BODY
else
    BODY
fi
\end{lstlisting}
\end{frame}

\begin{frame}[fragile]
\frametitle{\textbf{if} statements: exit code}
\begin{lstlisting}[language=bash]
# if-prog.sh
if mkdir temp; then
    # write to file only if mkdir emits exit code 0
    echo "temporary dir created" >> temp/log
fi
\end{lstlisting}
\end{frame}

\begin{frame}[fragile]
\frametitle{\textbf{if} statements: compare strings}
\begin{lstlisting}[language=bash]
# if-comp-str.sh
str=foo
if [[ "$str" = foo ]]; then
    echo "str is foo"
fi
\end{lstlisting}
Operators other than \cverb|=| include \cverb|!=|, \cverb|>|, and \cverb|<| with
their usual meaning.
\begin{block}{Note}
    \cverb|[[| is a command. In the above example, \cverb|"$str"|,
    \cverb|=|, \cverb|foo|, and \cverb|]]| are passed in as arguments to
    \cverb|[[|. Therefore, spaces after \cverb|[[| and before \cverb|]]| are
    \textbf{required}. \cverb|[[| has exit code just like any other command.
\end{block}
\end{frame}

\begin{frame}[fragile]
\frametitle{\textbf{if} statements: compare integers}
\begin{lstlisting}[language=bash]
# if-comp-int.sh
i=3
if [[ $i -eq 3 ]]; then
    echo "i is 3"
fi
\end{lstlisting}
Other operators:
\begin{table}
    \centering
    \begin{tabular}{ll}
        \cverb|-ge| & $\geq$ \\
        \cverb|-gt| & $>$    \\
        \cverb|-le| & $\leq$ \\
        \cverb|-lt| & $<$    \\
        \cverb|-ne| & $\neq$
    \end{tabular}
\end{table}
\end{frame}

\begin{frame}[fragile]
\frametitle{\textbf{if} statements: test strings}
\begin{lstlisting}[language=bash]
# if-test-str.sh
str="something"
if [[ -z $str ]]; then
    echo "str is empty"
fi
\end{lstlisting}
There is also \cverb|-n| for testing whether a string is not empty.
\end{frame}

\begin{frame}[fragile]
\frametitle{\textbf{if} statements: test files}
\begin{lstlisting}[language=bash]
# if-test-file.sh
file="something.txt"
if [[ -e $file ]]; then
    echo "$file exists"
fi
\end{lstlisting}
Other tests:
\begin{table}
    \centering
    \begin{tabular}{ll}
        \cverb|-f| & is a regular file \\
        \cverb|-d| & is a directory \\
    \end{tabular}
\end{table}
\end{frame}

\begin{frame}[fragile]
\frametitle{\textbf{if} statements: arithmetic}
\begin{lstlisting}[language=bash]
# if-arith.sh
i=3
if (( $i % 2 == 1 )); then
    echo "i is odd"
fi
\end{lstlisting}
\begin{block}{Note}
    \cverb|(( ))| is called \textbf{arithmetic evalution}, within which we can
    perform basic arithmetic operations like addition, division, etc. We use
    \cverb|==|, \cverb|!=|, \cverb|>|, and \cverb|<| for comparison. This is
    different from \cverb|[[ ]]|.
\end{block}
\end{frame}

\begin{frame}[fragile]
\frametitle{\textbf{for} statements}
The \tt{for} statement is usually used to:
\begin{itemize}
    \item Iterate over a range
    \item Iterate over a list of strings
\end{itemize}
It looks like:
\begin{lstlisting}[language=bash]
for VAR in LIST; do
    BODY
done
\end{lstlisting}
\end{frame}

\begin{frame}[fragile]
\frametitle{\textbf{for} statements: range}
\begin{lstlisting}[language=bash]
# for-timer.sh
for t in {10..1}; do
    echo "$t seconds left"
    sleep 1
done
\end{lstlisting}
\end{frame}

\begin{frame}[fragile]
\frametitle{\textbf{for} statements: globs}
\begin{lstlisting}[language=bash]
# for-backup.sh
for file in *; do
    cp $file $file.backup
done
\end{lstlisting}
\end{frame}

\begin{frame}[fragile]
\frametitle{\textbf{for} statements: list of words}
The previous two versions of for loops are not fundamentally different. The
range \cverb|{10..1}| expands to \cverb|10 9 8 7 6 5 4 3 2 1|, and \cverb|*|
expands to a list of all files in the current directory. A minimal example of
for loop is:
\begin{lstlisting}[language=bash]
# for-minimal.sh
for name in alice bob charlie; do
    echo "My name is $name."
done
\end{lstlisting}
\end{frame}

\begin{frame}[fragile]
\frametitle{\textbf{while} statements}
\cverb|while| statements looks like:
\begin{lstlisting}[language=bash]
while COMMAND; do
    BODY
done
\end{lstlisting}
\cverb|BODY| is repeated as long as \cverb|COMMAND| returns exit code 0.
\end{frame}

\begin{frame}[fragile]
\frametitle{Your turn}
Print the odd numbered lines of \tt{05-scripting/logos.txt}.
\pause
\begin{block}{Solution}
\begin{lstlisting}[language=bash]
IFS=$'\n'
count=0
for line in $(cat logos.txt); do
    count=$(( count + 1 ))
    if (( count % 2 == 1 )); then
        echo $line
    fi
done
\end{lstlisting}
\end{block}
\end{frame}

\begin{frame}[fragile]
\frametitle{The Ultimate Challenge}
Each line in \tt{05-scripting/obf.txt} begins with an 8-digit number.
Print every line with its number greater than any of the lines above.
\pause
\begin{block}{Solution}
\begin{lstlisting}[language=bash]
IFS=$'\n'
max=0
for line in $(cat obf.txt); do
    num=$(echo $line | grep -oE '^[0-9]{8}')
    if [[ $num -gt $max ]]; then
        echo $line
        max=$num
    fi
done
\end{lstlisting}
\end{block}
\end{frame}

\begin{frame}[fragile]
\frametitle{Environment Setup: \texttt{\textasciitilde/.bashrc}}
\texttt{\textasciitilde/.bashrc} is a script that runs automatically every time you open a new interactive terminal. 
\newline \newline
It is the perfect place to set environment variables and \textbf{aliases} (custom shortcuts).
\begin{lstlisting}[language=bash]
# Edit your ~/.bashrc
$ nano ~/.bashrc

# Add to the end of the file:
export PATH="$HOME/my_scripts:$PATH"
alias ll='ls -alF'
alias c='clear'
alias rm='rm -i' # Prompt before overwrite/delete
\end{lstlisting}
\begin{block}{Note}
After editing, you must run \ctexttt{source \textasciitilde/.bashrc} for changes to take effect in your current session!
\end{block}
\end{frame}

\begin{frame}[fragile]
\frametitle{Extension: Process Management}
When a program freezes, you need to manage its process.
\begin{itemize}
    \item \cverb|ps aux|: List all running processes.
    \item \cverb|pgrep [name]|: Find the Process ID (PID) by name.
    \item \cverb|kill [PID]|: Terminate a process by its PID.
    \item \cverb|pkill [name]|: Forcefully terminate processes by name.
\end{itemize}

\begin{lstlisting}[language=bash]
# Find the PID of python
$ pgrep python
12345

# Kill it violently by name (very useful!)
$ pkill -9 python
\end{lstlisting}
\end{frame}

\begin{frame}[fragile]
\frametitle{Extension: System Management}
Basic commands to control the machine's power state (requires \cverb|sudo|):
\begin{itemize}
    \item \cverb|shutdown|: Turn off the system.
    \item \cverb|reboot|: Restart the system immediately.
\end{itemize}

\begin{lstlisting}[language=bash]
$ sudo reboot          # Restart now

$ sudo shutdown now    # Shut down immediately
$ sudo shutdown +10    # Shut down in 10 minutes
$ sudo shutdown -c     # Cancel a pending shutdown
\end{lstlisting}
\end{frame}

\begin{frame}
\frametitle{Conclusion}
\begin{itemize}
    \item Regex matches apparent patterns in strings
    \item Many tools support regex, but beware of standards
    \item Bash scripts are for basic automation
    \item When unwieldy, consider alternatives
\end{itemize}
\end{frame}
